\documentclass[11pt,a4paper]{article}
\usepackage[margin=2.3cm]{geometry}
\usepackage{amsmath,amssymb,amsthm,bm,booktabs,array,hyperref,enumitem}
\usepackage{microtype}
\usepackage[dvipsnames]{xcolor}
\hypersetup{colorlinks=true,linkcolor=blue!50!black,citecolor=blue!50!black,urlcolor=blue!50!black}

\newtheorem{theorem}{Theorem}
\newtheorem{proposition}[theorem]{Proposition}
\newtheorem{lemma}[theorem]{Lemma}
\theoremstyle{remark}
\newtheorem{remark}[theorem]{Remark}
\theoremstyle{definition}
\newtheorem{assumption}{Assumption}

\newcommand{\diag}{\operatorname{diag}}
\newcommand{\dz}{\operatorname{dz}}
\newcommand{\sat}{\operatorname{sat}}
\newcommand{\clip}{\operatorname{clip}}
\newcommand{\sign}{\operatorname{sign}}
\newcommand{\Real}{\mathbb{R}}
\newcommand{\T}{^{\top}}
\newcommand{\test}[1]{\par\smallskip{\raggedright\renewcommand{\_}{\textunderscore\allowbreak}\noindent\textit{Numerical check:} \texttt{#1}.\par}}

\title{\textbf{MERLIN: Derivations and Proofs\\for the Bilateral Controller}}
\author{Companion document to the MERLIN paper and repository}
\date{}

\begin{document}
\maketitle

\begin{abstract}\noindent
This document states, for every equation of the revised MERLIN controller, the assumptions under which it
holds, the result, and a proof. Each statement that can be checked numerically is tied to a unit test in
\texttt{ros\_ws/src/controller/test/test\_derivations.py}. Section~\ref{sec:limits} lists what is \emph{not} proven.
It also explains, equation by equation, how the revised formulas differ from the equations of the submitted draft.
\end{abstract}

\tableofcontents

%% =====================================================================
\section{Scope, notation and standing assumptions}\label{sec:scope}
%% =====================================================================
The teleoperation pair is a ReactorX\,200 master and a ViperX\,300s slave with the forearm-roll joint locked at zero.
Each arm has $n=5$ active joints $q=[q_{\text{waist}},q_{\text{shoulder}},q_{\text{elbow}},q_{\text{wrist angle}},q_{\text{wrist rotate}}]^\top$.
The controller runs at $50$\,Hz, $\Delta t=0.02$\,s. Subscripts $m$ and $s$ denote master and slave. For a vector $x$, $\diag(x)$ is the
diagonal matrix, $\odot$ the elementwise product, $\sat_a(x)=\clip(x,-a,a)$ and $\dz_\delta(x)=\sign(x)\max(|x|-\delta,0)$ act
elementwise. $\mathcal F_\alpha\{x\}[k]=\alpha x[k]+(1-\alpha)\mathcal F_\alpha\{x\}[k-1]$ is a first-order IIR filter.

\begin{assumption}[rigid-body model]\label{as:rb}
For each arm, $M(q)=M(q)^\top\succ0$ is the joint-space inertia matrix, $C(q,\dot q)$ is the Coriolis matrix built from the
Christoffel symbols, and $G(q)$ the gravity torque, so that
\begin{equation}\label{eq:dyn}
M(q)\ddot q+C(q,\dot q)\dot q+G(q)=\tau_m+\tau_{ext}.
\end{equation}
$\tau_m$ is the motor torque and $\tau_{ext}$ the \emph{external torque}: everything not in the model, i.e.\ contact and
operator forces, joint friction and model error. The Christoffel construction gives $\dot M=C+C^\top$, hence
$\dot q^\top(\dot M-2C)\dot q=0$.
\end{assumption}

\begin{assumption}[actuation and sensing]\label{as:act}
Motor torque is $\tau_m=n_\mu\odot k_t\odot10^{-3}I$, with $I$ the present current in mA (as published by the Interbotix
\texttt{xs\_sdk}), $k_t$ the torque constant in N\,m/A, and $n_\mu$ the number of motors per joint (two for the
shoulder of the RX200 and for the shoulder and elbow of the VX300s). The current loop is much faster than $50$\,Hz.
Joint positions are measured; $\dot q$ is obtained from them.
\end{assumption}

\begin{assumption}[slave position loop]\label{as:loop}
Each slave joint is driven by a Dynamixel position loop, i.e.\ it accepts a position command $u$, not a torque. Seen from $u$,
joint $i$ behaves as a decoupled second-order system with a lumped, slowly varying acceleration disturbance $c_i$
(gravity sag, friction):
\begin{equation}\label{eq:loop}
\ddot q_i=\omega_i^{2}(u_i-q_i)-2\zeta_i\omega_i\dot q_i+c_i .
\end{equation}
$\omega_i,\zeta_i$ are identified per joint (\texttt{tools/identify\_position\_loop.py}). This is a modelling assumption, validated
in simulation only.
\end{assumption}

\paragraph{Sign convention.}
$\tau_{ext}$ is the torque the environment applies \emph{on} the joint. The torque the controller renders on the operator
points in the same direction as the environment's torque on the slave (the operator is pushed back).

%% =====================================================================
\section{External-torque estimation}\label{sec:obs}
%% =====================================================================
\subsection{The static residual and why it fails during motion}
The draft estimates $\hat\tau_{ext}=\tau_m-(C-G)$; with Assumption~\ref{as:rb} the correct static residual is the following.

\begin{lemma}[static residual]\label{lem:static}
Define $\hat\tau^{\text{static}}:=G(q)+C(q,\dot q)\dot q-\tau_m$. Then $\hat\tau^{\text{static}}=\tau_{ext}-M(q)\ddot q$.
\end{lemma}
\begin{proof}
Solve \eqref{eq:dyn} for $\tau_{ext}$: $\tau_{ext}=M\ddot q+C\dot q+G-\tau_m$, and subtract $M\ddot q$.
\end{proof}
So the residual is correct only when $M\ddot q$ is negligible. On a position-controlled arm the position loop accelerates the
links, and $M\ddot q$ appears as a false contact torque. In the MuJoCo benchmark the shoulder inertia is
$M_{22}\approx0.23$--$0.30$\,kg\,m$^2$ and the free-motion false torque grows with speed as $0.25,\,0.48,\,1.08$\,N\,m
at $1\times,2\times,3\times$ the motion speed; its correlation with $-M\ddot q$ is $0.67,\,0.81,\,0.92$.
At low speed friction, mass mismatch and current noise dominate; at high speed $M\ddot q$ does.
\test{test\_static\_residual\_equals\_tau\_ext\_minus\_Mqdd, test\_inverse\_dynamics\_decomposition}

\subsection{Generalised-momentum observer}
Define the generalised momentum $p:=M(q)\dot q$ and $\beta(q,\dot q):=G(q)-C(q,\dot q)^\top\dot q$.

\begin{lemma}[momentum balance]\label{lem:mom}
Under Assumption~\ref{as:rb}, $\dot p=\tau_m+\tau_{ext}-\beta(q,\dot q)$.
\end{lemma}
\begin{proof}
$\dot p=\dot M\dot q+M\ddot q$. Substituting \eqref{eq:dyn}, $M\ddot q=\tau_m+\tau_{ext}-C\dot q-G$, and $\dot M=C+C^\top$,
$\dot p=(C+C^\top)\dot q+\tau_m+\tau_{ext}-C\dot q-G=\tau_m+\tau_{ext}+C^\top\dot q-G$.
\end{proof}

\begin{theorem}[observer]\label{thm:obs}
Let $K_O=\diag(k_1,\dots,k_n)$, $k_i>0$, and define
\begin{equation}\label{eq:obs}
r(t)=K_O\Big(p(t)-p(0)-\int_0^t\big(\tau_m+C^\top\dot q-G+r\big)\,ds\Big).
\end{equation}
Then $r(0)=0$ and $\dot r=K_O(\tau_{ext}-r)$. In particular $r_i$ is a first-order low-pass filter of $\tau_{ext,i}$ with cutoff
$k_i$\,rad/s, $r_i(t)=\int_0^te^{-k_i(t-s)}k_i\,\tau_{ext,i}(s)\,ds$, and $r\to\tau_{ext}$ for a constant $\tau_{ext}$.
No acceleration is needed.
\end{theorem}
\begin{proof}
$r(0)=0$ by construction. Differentiating \eqref{eq:obs}, $\dot r=K_O(\dot p-\tau_m-C^\top\dot q+G-r)$. By Lemma~\ref{lem:mom},
$\dot p-\tau_m-C^\top\dot q+G=\tau_{ext}$. The solution of the scalar linear ODE $\dot r_i=k_i(\tau_{ext,i}-r_i)$, $r_i(0)=0$ is the stated
convolution; for $\tau_{ext}$ constant it equals $\tau_{ext,i}(1-e^{-k_it})$.
\end{proof}
\test{test\_Mdot\_equals\_C\_plus\_Ct, test\_observer\_converges\_to\_constant\_external\_torque}

\begin{remark}[bias from model error]\label{rem:bias}
If the implementation uses imperfect $\tilde M,\tilde G,\tilde C,\tilde k_t$, repeating the proof gives
$\dot{\tilde r}=K_O(\tau_{ext}+\delta-\tilde r)$ with
\[
\delta=(\tau_m-\tilde\tau_m)+(\tilde\beta-\beta)+\Delta\dot M\,\dot q+\Delta M\,\ddot q,\qquad\Delta M=\tilde M-M.
\]
So the estimate is the low-passed external torque \emph{plus} a bias from the torque constants, from gravity and Coriolis errors, and from the
inertia error times the acceleration (which is small: $\Delta M\ddot q$, not $M\ddot q$). In the simulation, link masses are $8\%$ higher than the
controller's model, so the gravity term alone gives $\delta\approx0.08\,G\approx0.2$\,N\,m at the shoulder and elbow
($G=[-2.3,-3.0]$\,N\,m at the reference pose), the same order as the $0.25$--$0.3$\,N\,m RMS error measured in the estimation study. On the real arms,
$k_t$ and the link masses must be calibrated with known weights (\texttt{docs/HARDWARE\_VALIDATION.md}).
\end{remark}

\begin{proposition}[discrete observer]\label{prop:disc}
Implement \eqref{eq:obs} at sample time $\Delta t$ with the rectangle rule: $S[k]=S[k-1]+\Delta t\,(\tau_m[k]-\beta[k]+r[k-1])$,
$r[k]=K_O(p[k]-p[0]-S[k])$. If $p[k]-p[k-1]=\Delta t(\tau_m[k]-\beta[k]+\tau_{ext}[k])+\varepsilon[k]$, with $\varepsilon[k]$ the
quadrature error of the integral of $\dot p$, then
\begin{equation}\label{eq:iir}
r[k]=(1-K_O\Delta t)\,r[k-1]+K_O\Delta t\,\tau_{ext}[k]+K_O\varepsilon[k].
\end{equation}
It is a stable first-order IIR filter iff $0<k_i\Delta t<2$, and monotone iff $k_i\Delta t<1$. With $K_O=15$\,rad/s,
$\Delta t=0.02$\,s: $a=K_O\Delta t=0.3$, equivalent time constant $-\Delta t/\ln(1-a)=56$\,ms
(the continuous value is $1/K_O=67$\,ms).
\end{proposition}
\begin{proof}
Subtract the update for $r[k-1]$ from the one for $r[k]$:
$r[k]-r[k-1]=K_O\big((p[k]-p[k-1])-\Delta t(\tau_m[k]-\beta[k]+r[k-1])\big)$. Insert the balance for $p[k]-p[k-1]$ to get
$r[k]-r[k-1]=K_O\Delta t(\tau_{ext}[k]-r[k-1])+K_O\varepsilon[k]$, which is \eqref{eq:iir}. The pole is $1-k_i\Delta t$, inside the unit
circle iff $0<k_i\Delta t<2$.
\end{proof}
\test{test\_discrete\_observer\_recursion\_is\_exact\_iir (the recursion holds to $10^{-10}$ when $\varepsilon=0$)}

\begin{remark}[unfiltered inputs; empirical]
The implementation feeds the observer the unfiltered velocity and torque. The velocity IIR filter of the draft ($\alpha_v=0.3$) has a mean delay of
$(1-\alpha_v)/\alpha_v=2.3$ samples $=47$\,ms, comparable with the observer time constant. Feeding the filtered velocity delayed $p=M\dot q$
relative to $\int\tau$ and, in the simulation, produced a master--slave limit cycle (free-motion RMSE $0.113$\,rad, measured with the first solver version, against $0.007$\,rad with the observer fed raw velocity).
This is an empirical observation, not a theorem.
\end{remark}

%% =====================================================================
\section{Slave admittance: exact discretisation}\label{sec:adm}
%% =====================================================================
For each slave joint, the compliance offset $q_e$ obeys the desired impedance
\begin{equation}\label{eq:adm}
M_d\ddot q_e+D_d\dot q_e+K_dq_e=u,\qquad u=\dz_\delta(\hat\tau_{ext,s}),\qquad q_{s,\mathrm{ref}}=Sq_m+o+q_e,
\end{equation}
with design parameters $M_d,D_d,K_d>0$ (not the physical inertia). With $x=[q_e,\dot q_e]^\top$,
$A=\begin{bmatrix}0&1\\-K_d/M_d&-D_d/M_d\end{bmatrix}$, $B=\begin{bmatrix}0\\1/M_d\end{bmatrix}$ we have $\dot x=Ax+Bu$.

\begin{theorem}[exact zero-order-hold discretisation]\label{thm:zoh}
If $u$ is constant on $[k\Delta t,(k+1)\Delta t)$, then $x[k+1]=A_dx[k]+B_du[k]$ exactly, with $A_d=e^{A\Delta t}$,
$B_d=\int_0^{\Delta t}e^{As}B\,ds$, and
\begin{equation}\label{eq:vanloan}
\exp\!\Big(\Delta t\begin{bmatrix}A&B\\0&0\end{bmatrix}\Big)=\begin{bmatrix}A_d&B_d\\0&1\end{bmatrix}.
\end{equation}
\end{theorem}
\begin{proof}
Variation of constants gives $x(t_{k+1})=e^{A\Delta t}x(t_k)+\int_0^{\Delta t}e^{A(\Delta t-s)}B\,u[k]\,ds$, and the substitution $s\mapsto\Delta t-s$ turns the integral into
$B_du[k]$. For \eqref{eq:vanloan}, let $E(t)=\exp(t\bar A)$ with $\bar A=\begin{bmatrix}A&B\\0&0\end{bmatrix}$ and write $E(t)=\begin{bmatrix}E_{11}&F\\0&1\end{bmatrix}$.
From $\dot E=\bar AE$: $\dot E_{11}=AE_{11}$, $E_{11}(0)=I$, so $E_{11}=e^{At}$; and $\dot F=AF+B$, $F(0)=0$, so $F=\int_0^te^{A(t-s)}B\,ds=\int_0^te^{As}B\,ds$
(Van~Loan's formula \cite{vanloan1978}).
\end{proof}
\test{test\_zoh\_matches\_fine\_integration (agrees with a $10^{-12}$-tolerance ODE solve to $10^{-10}$)}

\begin{proposition}[stability; failure of explicit Euler]\label{prop:euler}
(i) For $M_d,D_d,K_d>0$ every eigenvalue $\lambda$ of $A$ has $\operatorname{Re}\lambda<0$, so the eigenvalues $e^{\lambda\Delta t}$ of $A_d$ satisfy
$|e^{\lambda\Delta t}|=e^{\operatorname{Re}\lambda\,\Delta t}<1$ for \emph{every} $\Delta t>0$. (ii) The static deflection is $q_e=u/K_d$.
(iii) Explicit Euler, $x[k+1]=(I+\Delta tA)x[k]+\Delta tBu[k]$, is stable iff $|1+\lambda\Delta t|<1$ for both eigenvalues.
\end{proposition}
\begin{proof}
(i) The eigenvalues are the roots of $M_ds^2+D_ds+K_d=0$; a real quadratic with all coefficients of the same sign has both roots in the open left half plane
(Routh--Hurwitz). (ii) Set $\dot x=0$. (iii) The eigenvalues of $I+\Delta tA$ are $1+\lambda\Delta t$.
\end{proof}

\begin{remark}[worked example: the original wrist]
The original code used the robot inertia as $M_d$: for the wrist rotate, $M_d=7.29\times10^{-4}$\,kg\,m$^2$, with $D_d=0.4$, $K_d=10$.
The roots are $\lambda=-26.3$ and $-522.4$\,s$^{-1}$ ($\omega_n=117$\,rad/s, $\zeta=2.34$). At $\Delta t=20$\,ms,
$|1+\lambda\Delta t|=0.47$ and $9.45$: Euler diverges, whereas $|e^{\lambda\Delta t}|=0.59$ and $3\times10^{-5}$. The original code kept the
state bounded only with ad-hoc clamps ($\pm0.05$\,rad, $\ddot q_e\le30$\,rad/s$^2$), which chatter. The revised design values
$M_d=[0.30,0.40,0.30,0.10,0.04]$, $D_d=[6,8,6,2,0.8]$, $K_d=[30,40,30,10,4]$ give $\omega_n=\sqrt{K_d/M_d}\approx10$\,rad/s, $\zeta=1$.
\end{remark}
\test{test\_euler\_unstable\_for\_original\_wrist\_admittance, test\_admittance\_static\_deflection\_is\_tau\_over\_K}

\begin{remark}
The implementation saturates $|q_e|\le q_{e,\max}$ and zeroes $\dot q_e$ when pushing into the limit (anti-windup). The saturated filter is bounded
by construction; Theorem~\ref{thm:zoh} and Proposition~\ref{prop:euler} describe the unsaturated range.
\end{remark}

%% =====================================================================
\section{Master law}\label{sec:master}
%% =====================================================================
The RX200 master runs in Dynamixel current mode, so the commanded current is applied as torque. Its \emph{measured} current is therefore, up to noise and
delay, the \emph{commanded} current.

\begin{lemma}[the original law integrates its own command]\label{lem:pole}
Let the measured torque be the command delayed by one sample and IIR-filtered with coefficient $\alpha$,
$\tilde\tau[k]=\mathcal F_\alpha\{\tau_{cmd}[k-1]\}$. The original law is
$\tau_{cmd}=\hat G-D_{md}\dot q+\hat\tau_{ext,m}+\tau_f$ with $\hat\tau_{ext,m}=\tilde\tau-C-\hat G$ (its own estimator).
Then $\tau_{cmd}[k]=\tilde\tau[k]+b[k]$ with $b=\tau_f-D_{md}\dot q-C\dot q$, which has a pole at $z=1$.
A constant $b\ne0$ makes $\tau_{cmd}$ ramp by $\alpha b$ per sample until the clamp.
\end{lemma}
\begin{proof}
Substitute $\hat\tau_{ext,m}$: the $\hat G$ terms cancel, so $\tau_{cmd}[k]=\tilde\tau[k]+b[k]$. With $\tilde\tau[k]=\alpha\tau_{cmd}[k-1]+(1-\alpha)\tilde\tau[k-1]$
the transfer function from $b$ to $\tau_{cmd}$ is $\big(1-(1-\alpha)z^{-1}\big)/(1-z^{-1})$: its DC gain of the feedback path is $1$, so the pole is at $z=1$.
For constant $b$: $\tilde\tau[k]=\tilde\tau[k-1]+\alpha b$ (induction from $\tilde\tau[k]=\alpha(\tilde\tau[k-1]+b)+(1-\alpha)\tilde\tau[k-1]$).
\end{proof}
\test{test\_original\_master\_law\_has\_a\_pole\_at\_one}
Any constant bias in $b$ (a gravity-model error, a sensor offset, the slave's coupling term) is integrated until the torque clamp: this is the ``rough or no change'' behaviour reported
on hardware. The revised law removes $\hat\tau_{ext,m}$:
\begin{equation}\label{eq:master}
\tau_{m,cmd}=\hat G_m(q_m)-D_{md}\dot q_m+\gamma\,\tau_f ,
\end{equation}
where $\tau_f$ depends on slave data and on $q_m-q_s$ only, and $\gamma\in[0,1]$ is the tank gate of Section~\ref{sec:tank}.
There is no path from the master's measured current to its command, so the pole of Lemma~\ref{lem:pole} does not arise.

\begin{proposition}[passivity of the master with the revised law]\label{prop:masterpass}
Let $\hat G_m=G_m$ and let the operator apply $\tau_h=J_m^\top F_h$. With $V=\tfrac12\dot q_m^\top M_m\dot q_m$,
\[
\dot V=\dot q_m^\top\tau_h-\dot q_m^\top D_{md}\dot q_m+\gamma\,\dot q_m^\top\tau_f .
\]
\end{proposition}
\begin{proof}
$\dot V=\dot q^\top M\ddot q+\tfrac12\dot q^\top\dot M\dot q=\dot q^\top(\tau_{m,cmd}+\tau_h-C\dot q-G)+\tfrac12\dot q^\top\dot M\dot q$. By Assumption~\ref{as:rb},
$\dot q^\top(\tfrac12\dot M-C)\dot q=0$. Insert \eqref{eq:master}; the gravity terms cancel.
\end{proof}
With $\tau_f=0$ the master is output-strictly passive from $\tau_h$ to $\dot q_m$. The last term, $\gamma\dot q_m^\top\tau_f=\gamma P_{out}$, is the power the force
channel delivers to the operator. It is exactly the quantity the energy tank of Section~\ref{sec:tank} bounds. Gravity compensation is approximate on hardware, so the statement holds up to
$\dot q^\top(\hat G-G)$.

%% =====================================================================
\section{Force reflection between kinematically different arms}\label{sec:map}
%% =====================================================================
Reflecting joint torque one-to-one renders the wrong Cartesian force when the arms differ. At the reference pose the RX200 reaches $0.388$\,m
and the VX300s $0.530$\,m (ratio $1.37$).

\begin{proposition}[damped least-squares wrench]\label{prop:dls}
Let $J_s\in\Real^{6\times n}$ be the slave end-effector Jacobian (full column rank $n=5$) and $\lambda>0$. The wrench
\begin{equation}\label{eq:dls}
\hat F_s=\arg\min_F\ \|J_s^\top F-\hat\tau_{ext,s}\|^2+\lambda^2\|F\|^2=J_s\big(J_s^\top J_s+\lambda^2I\big)^{-1}\hat\tau_{ext,s}
\end{equation}
satisfies $J_s^\top\hat F_s=\hat\tau_{ext,s}-\lambda^2(J_s^\top J_s+\lambda^2I)^{-1}\hat\tau_{ext,s}$, i.e.\ $J_s^\top\hat F_s\to\hat\tau_{ext,s}$ as $\lambda\to0$.
\end{proposition}
\begin{proof}
Setting the gradient to zero, $J_s(J_s^\top F-\hat\tau)+\lambda^2F=0$, so $(J_sJ_s^\top+\lambda^2I)F=J_s\hat\tau$. The push-through identity
$(J_sJ_s^\top+\lambda^2I)^{-1}J_s=J_s(J_s^\top J_s+\lambda^2I)^{-1}$ gives \eqref{eq:dls}. Then $J_s^\top\hat F_s=J_s^\top J_s(J_s^\top J_s+\lambda^2I)^{-1}\hat\tau=
\hat\tau-\lambda^2(J_s^\top J_s+\lambda^2I)^{-1}\hat\tau$.
\end{proof}
The $6$-dimensional wrench is not unique because five joints cannot produce one wrench direction; (\ref{eq:dls}) selects the minimum-norm one.
The torque rendered on the master is
\begin{equation}\label{eq:force}
\tau_f=k_{of}\,J_m^\top S_f\,\hat F_s-k_{om}K_{md}(Sq_m+o-q_s),\qquad S_f=\diag(s_f I_3,\,s_\mu I_3).
\end{equation}
(the second term is the position channel; $\hat\tau_{ext,s}$ in \eqref{eq:dls} is first passed through $\dz_\delta$).

\begin{proposition}[consistency]\label{prop:consistent}
If the two arms have the same Jacobian at their configurations, $S_f=I$ and $\lambda\to0$, then $J_m^\top\hat F_s=\hat\tau_{ext,s}$: the Jacobian map reduces to the identity map.
For two different arms and the same joint torque $\tau$, the identity map yields a master-side force $F_{\rm id}=(J_m^\top)^{+}\tau$ whereas the force that produced $\tau$ on the slave is
$F_s=(J_s^\top)^{+}\tau$; at the reference pose $\|F_{\rm id}\|/\|F_s\|\approx1.3$.
\end{proposition}
\begin{proof}
The first claim is Proposition~\ref{prop:dls} with $J_m=J_s$. The ratio follows from $\tau=J^\top F$ (virtual work, $\delta W=\tau^\top\delta q=F^\top J\delta q$): a shorter arm needs a larger force for the same joint torque.
\end{proof}
\test{test\_dls\_wrench\_reproduces\_joint\_torque, test\_identical\_arms\_give\_identity\_map, test\_different\_arms\_over\_render\_joint\_torque\_reflection}
The simulation shows the effect: reflecting the vertical force with the identity map over-renders it by a factor of $1.40$, the Jacobian map by $1.10$. Of the remaining excess,
the position-channel term adds force when the slave sags under a sustained load.

\begin{remark}[power]
$\Gamma=J_m^\top S_fJ_s^{+\top}$ is not power-preserving between the arms, because the arms are coupled in joint space ($q_m\leftrightarrow q_s$), not in task space. The slave-side power
$\hat\tau^\top\dot q_s=\hat F^\top\dot x_s$ and the master-side power $\tau_f^\top\dot q_m=\hat F^\top S_f^\top\dot x_m$ are equal only if $\dot x_m=S_f^{-\top}\dot x_s$, which joint-space mapping does not enforce.
This is why the force channel is wrapped by the energy tank.
\end{remark}

%% =====================================================================
\section{Slave MPC on the position loop}\label{sec:mpc}
%% =====================================================================
\subsection{What the draft's Eq.~(7) said, and what was implemented}
The draft's MPC minimises over torques $\tau_{0:N-1}$ of a rigid-body model and applies $\tau_0^\star$. The VX300s accepts position commands; no torque was ever sent to it.
In the code the decision vector was $[dq_{0:N-1},\tau_{0:N-1}]$ but $dq$ never entered the rollout: the cost is $dq^\top Rdq+\dots$, so its minimiser is $dq=0$ for any reference, and the value used, $dq_0^\star$, was identically zero.
The slave moved only through the smoothing step $q\leftarrow0.95q+0.05q_{ref}$, a first-order lag of $0.39$\,s at $50$\,Hz. The joint-limit constraint of the draft was not enforced either
(\texttt{q\_min}, \texttt{q\_max} were stored and never used). The revised formulation below replaces all of this.

\subsection{Model and prediction}
By Theorem~\ref{thm:zoh} with $A=\begin{bmatrix}0&1\\-\omega^2&-2\zeta\omega\end{bmatrix}$, $B=\begin{bmatrix}0\\\omega^2\end{bmatrix}$, joint $i$ of \eqref{eq:loop} gives $x[k+1]=A_dx[k]+B_du[k]$ (the disturbance $c$ is handled in Section~\ref{sec:offset}).
For one joint, with horizon $N$ and model state $\hat x_0$, the predicted positions and velocities are
\begin{equation}\label{eq:pred}
\hat q_k=e_1^\top A_d^{k}\hat x_0+\sum_{j=0}^{k-1}e_1^\top A_d^{k-1-j}B_d\,u_j+\hat d,\qquad
\hat v_k=e_2^\top A_d^{k}\hat x_0+\sum_{j=0}^{k-1}e_2^\top A_d^{k-1-j}B_d\,u_j,\qquad k=1,\dots,N,
\end{equation}
that is $\hat q=\Phi_q\hat x_0+G_qU+\hat d\mathbf 1$ and $\hat v=\Phi_v\hat x_0+G_vU$ with $U=[u_0,\dots,u_{N-1}]^\top$ and $G_q,G_v$ lower-triangular Toeplitz.

\subsection{The quadratic program}
Let $r_k$ be the reference preview ($r_k=q_{s,\mathrm{ref}}+k\Delta t\,S\dot q_m$), $w=(Q,\dots,Q,P)$, $W=\diag(w)$, $D$ the unit lower-bidiagonal difference matrix ($(DU)_k=u_k-u_{k-1}$) and $d_0=[u_{-1},0,\dots,0]^\top$ with $u_{-1}$ the last applied command:
\begin{equation}\label{eq:mpc}
\min_{U}\ \ (\hat q-r)^\top W(\hat q-r)+R\,\|DU-d_0\|^2\quad\text{s.t.}\quad
\begin{aligned}
&q_{\min}\le u_k\le q_{\max},\\
&|u_k-u_{k-1}|\le\Delta u_{\max},\\
&|\hat v_k|\le\dot q_{\max}.
\end{aligned}
\end{equation}

\begin{proposition}[condensed form]\label{prop:qp}
Let $a=\Phi_q\hat x_0+\hat d\mathbf 1-r$. Then \eqref{eq:mpc} equals $\tfrac12U^\top HU+f^\top U+\mathrm{const}$ subject to $l\le CU\le h$, with
\begin{equation}\label{eq:H}
H=2\big(G_q^\top WG_q+R\,D^\top D\big),\quad f=2G_q^\top Wa-2R\,D^\top d_0,\quad C=\begin{bmatrix}I\\D\\G_v\end{bmatrix},
\end{equation}
$l=[q_{\min}\mathbf1;\ d_0-\Delta u_{\max}\mathbf1;\ -\dot q_{\max}\mathbf1-\Phi_v\hat x_0]$ and $h$ analogously. $H\succ0$ for $R>0$, so a feasible problem has a unique minimiser.
\end{proposition}
\begin{proof}
Expand $(G_qU+a)^\top W(G_qU+a)+R(DU-d_0)^\top(DU-d_0)$ and collect the terms quadratic and linear in $U$. $D$ is unit lower bidiagonal, hence invertible, so $D^\top D\succ0$;
$G_q^\top WG_q\succeq0$ since $W\succeq0$. The bounds follow by writing the three constraint groups in the variable $CU$.
\end{proof}

\subsection{Solving it in about a millisecond: ADMM}\label{sec:admm}
With the auxiliary variable $Z=CU$, ADMM with proximal weight $\sigma>0$ and penalty $\rho>0$ iterates
\begin{align}
U^{+}&=(H+\sigma I+\rho C^\top C)^{-1}\big(\sigma U-f+C^\top(\rho Z-y)\big),\label{eq:admmU}\\
Z^{+}&=\Pi_{[l,h]}\big(CU^{+}+y/\rho\big),\qquad y^{+}=y+\rho\,(CU^{+}-Z^{+}),\label{eq:admmZ}
\end{align}
(the OSQP scheme without relaxation \cite{stellato2020osqp}, a special case of ADMM \cite{boyd2011admm}).
\begin{proposition}[ADMM]\label{prop:admm}
(i) \eqref{eq:admmU} is the minimiser of the augmented Lagrangian $\tfrac12U^\top HU+f^\top U+\tfrac\sigma2\|U-U^k\|^2+\tfrac\rho2\|CU-Z^k+y^k/\rho\|^2$.
(ii) A fixed point $(U,Z,y)$ of \eqref{eq:admmU}--\eqref{eq:admmZ} satisfies the KKT conditions of \eqref{eq:mpc}: $HU+f+C^\top y=0$, $Z=CU\in[l,h]$, $y\in N_{[l,h]}(Z)$.
(iii) The matrix $H+\sigma I+\rho C^\top C$ does not depend on $\hat x_0,r,u_{-1},\hat d$ or the bounds, so it is inverted once; an iteration costs one $N\times N$ product and one clipping.
(iv) For any $\rho,\sigma>0$ and a feasible convex QP the iterates converge to a KKT point \cite{boyd2011admm,stellato2020osqp}.
\end{proposition}
\begin{proof}
(i) The gradient condition is $HU+f+\sigma(U-U^k)+\rho C^\top(CU-Z^k+y^k/\rho)=0$, i.e.\ \eqref{eq:admmU}. (ii) At a fixed point the $y$-update gives $CU=Z$; the $Z$-update gives $Z=\Pi(Z+y/\rho)$, i.e.\ $y/\rho$ in the normal cone of
$[l,h]$ at $Z$; and \eqref{eq:admmU} with $U^+=U^k=U$, $Z^k=Z$ reduces to $HU+\rho C^\top CU=-f+\rho C^\top Z-C^\top y$, which with $CU=Z$ is $HU+f+C^\top y=0$. (iii) is read off. (iv) is the cited result.
\end{proof}
\paragraph{Choice of $\rho$ (a correction).}
A fixed number of iterations only approximates the optimum, and the approximation depends strongly on $\rho$. With $\rho=10$ (the first implementation) and 60 iterations, the first move was up to $26$\,mrad
from the exact solution from a cold start, and on average $13$\,mrad (maximum $48$\,mrad) when warm-started in closed loop: the controller was effectively a different, less regularised design than the documented QP.
With $\rho=0.3$ the warm-started first move is within $0.03$\,mrad on average ($0.2$\,mrad maximum) of the exact solution, at the same cost ($\approx0.9$\,ms for five joints). The experiments in the repository were re-run with $\rho=0.3$.
\test{test\_admm\_solves\_the\_stated\_qp (reference: SLSQP on the same QP data), test\_applied\_command\_always\_respects\_limits}

\begin{proposition}[the applied command is always admissible]\label{prop:feas}
After the solve, $u_0^\star$ is clipped to $[q_{\min},q_{\max}]$ and to $[u_{-1}-\Delta u_{\max},u_{-1}+\Delta u_{\max}]$. Since $u_{-1}$ satisfied the same limits, the intersection is non-empty,
and the applied command satisfies the position and rate limits regardless of solver accuracy (the predicted-velocity bound is enforced by the QP only).
\end{proposition}

\subsection{Offset-free tracking}\label{sec:offset}
The Dynamixel position loop has no gravity feed-forward, so under constant load joint $i$ settles at $u+c/\omega^2$ rather than $u$. The model \eqref{eq:pred} cannot predict this.
The controller runs the model in parallel with the plant (\emph{internal-model} form): $\hat x[k+1]=A_d\hat x[k]+B_du[k]$, and estimates an output disturbance
\begin{equation}\label{eq:dhat}
\hat d[k+1]=(1-\alpha_d)\hat d[k]+\alpha_d\big(q[k]-\hat x_1[k]\big),\qquad 0<\alpha_d\le1 .
\end{equation}
\begin{proposition}[offset-free]\label{prop:offsetfree}
Let the plant be \eqref{eq:loop} with constant $c$ and let $u\equiv u_\infty$. Then $\hat d\to c/\omega^2$. Moreover, if the reference is constant, $r\equiv r_\infty$, and $u_\infty=r_\infty-\hat d$ lies within the limits,
then $(u_\infty,\ \hat d=c/\omega^2)$ is an equilibrium of the closed loop at which the plant position equals $r_\infty$.
\end{proposition}
\begin{proof}
The plant converges to $q_\infty=u_\infty+c/\omega^2$ ($A_d$ is Schur by Theorem~\ref{thm:zoh}(i)); the parallel model converges to $\hat x_{1,\infty}=u_\infty$. Hence $q-\hat x_1\to c/\omega^2$ and the stable recursion \eqref{eq:dhat}
(pole $1-\alpha_d$) converges to the same value. For the second claim, with $\hat d=c/\omega^2$ and $U\equiv u_\infty=r_\infty-\hat d$ the predictions satisfy $\hat q_k+\hat d=r_\infty$ for all $k$ and $\Delta u=0$, so the cost of \eqref{eq:mpc} is $0$,
its global minimum; by Proposition~\ref{prop:qp} the minimiser is unique, so the MPC outputs $u_\infty$, and the plant sits at $u_\infty+c/\omega^2=r_\infty$.
\end{proof}
\test{test\_offset\_free\_estimate\_converges\_and\_removes\_steady\_state\_error (sag $48$\,mrad removed; $\hat d$ within $2\%$)}
This is the standard offset-free MPC mechanism \cite{muske2002,pannocchia2003}. Proposition~\ref{prop:offsetfree} establishes the existence of the zero-error equilibrium, not its asymptotic stability; the latter is checked in simulation only.

\subsection{What is not claimed}
No terminal cost or terminal set with a Lyapunov argument is used, so \emph{closed-loop stability of the MPC is not proven}; $N=10$ and the weights were chosen by simulation. The prediction model is decoupled per joint and validated in simulation only (Section~\ref{sec:limits}).

%% =====================================================================
\section{Energy tank on the force channel}\label{sec:tank}
%% =====================================================================
Let $P_{out}[k]=\tau_f[k]^\top\dot q_m[k]$ be the power the (unscaled) force channel delivers to the operator--master system, and
$P_{in}[k]=\eta\big(\dot q_m^\top D_{md}\dot q_m+\dot q_e^\top D_{sd}\dot q_e\big)[k]\ge0$ ($0<\eta\le1$) the power the controller dissipates through its two damping terms. The tank stores
$T[k]\in[T_{\min},T_{\max}]$, $0<T_{\min}\le T[0]\le T_{\max}$, and a gate $\gamma[k]\in[0,1]$ scales the force:
\begin{align}
\gamma[k]&=\begin{cases}1&\text{if }P_{out}[k]\le0\ \text{ or }\ b[k]\ge\Delta t\,P_{out}[k],\\[2pt]\dfrac{b[k]}{\Delta t\,P_{out}[k]}&\text{otherwise,}\end{cases}\qquad b[k]=T[k]+\Delta tP_{in}[k]-T_{\min},\label{eq:gamma}\\
T[k+1]&=\min\big\{T_{\max},\ T[k]+\Delta t\,(P_{in}[k]-\gamma[k]P_{out}[k])\big\}.\label{eq:tank}
\end{align}

\begin{theorem}[bounded delivered energy]\label{thm:tank}
For any signals $\tau_f,\dot q_m,P_{in}\ge0$ and every $K\ge1$:
(i) $\gamma[k]\in[0,1]$ and $T[k]\ge T_{\min}$ for all $k$;
(ii) the energy delivered by the gated force channel is bounded:
\begin{equation}\label{eq:bound}
\sum_{k=0}^{K-1}\Delta t\,\gamma[k]P_{out}[k]\ \le\ T[0]-T_{\min}+\sum_{k=0}^{K-1}\Delta t\,P_{in}[k].
\end{equation}
\end{theorem}
\begin{proof}
(i) Induction on $T[k]\ge T_{\min}$. Since $P_{in}\ge0$, $b[k]\ge0$, so $\gamma\ge0$; and $\gamma\le1$ by construction. If $P_{out}\le0$, $\gamma=1$ and
$T[k+1]=\min\{T_{\max},T[k]+\Delta t(P_{in}+|P_{out}|)\}\ge\min\{T_{\max},T[k]\}\ge T_{\min}$. If $P_{out}>0$ and $b\ge\Delta tP_{out}$, then $\gamma=1$ and $T[k]+\Delta t(P_{in}-P_{out})=b+T_{\min}-\Delta tP_{out}\ge T_{\min}$. Otherwise
$\gamma\Delta tP_{out}=b$ and $T[k]+\Delta tP_{in}-b=T_{\min}$. In every case the argument of the $\min$ in \eqref{eq:tank} is $\ge T_{\min}$ and $T_{\max}\ge T_{\min}$.
(ii) Write \eqref{eq:tank} as $T[k+1]=T[k]+\Delta t(P_{in}-\gamma P_{out})-c[k]$ with the clipping loss $c[k]\ge0$. Summing over $k$,
$\sum\Delta t\gamma P_{out}=T[0]-T[K]+\sum\Delta tP_{in}-\sum c[k]\le T[0]-T_{\min}+\sum\Delta tP_{in}$, using $T[K]\ge T_{\min}$ and $c\ge0$.
\end{proof}
\test{test\_tank\_energy\_bound\_theorem (random signals, 5000 steps), test\_energy\_tank\_never\_below\_minimum}

\paragraph{Reading of the result.}
Whatever the force estimate does, the force channel can return to the operator at most the energy that the controller has dissipated plus the initial reserve $T[0]-T_{\min}$; $T_{\max}$ caps how much can be released in a burst.
Combined with Proposition~\ref{prop:masterpass}, the master with the gated feedback has $\dot V\le\dot q^\top\tau_h-\dot q^\top D\dot q+\gamma P_{out}$, with the last term energetically bounded by \eqref{eq:bound}.

\paragraph{What it does not say.}
\begin{itemize}[leftmargin=*,itemsep=1pt]
\item \eqref{eq:bound} uses the Riemann sum $\Delta t\,\tau_f^\top\dot q_m$ for the energy of a sample; it is exact only to the discretisation error of that product. An exact zero-order-hold energy would use $\tau_f[k]^\top(q_m[k+1]-q_m[k])$, which needs the next sample.
\item It bounds the energy delivered by the \emph{master-side controller}. With communication delay the slave-side quantities in $\tau_f$ are stale, and a delayed channel is not passive in general \cite{anderson1989}; the tank does not address this.
The simulations confirm it: the tank enlarges the stable force gain without delay (from $k_{of}\approx0.6$ to $\approx1.0$--$1.5$) and gives no benefit at $20$--$60$\,ms of delay. A two-port tank or time-domain passivity control with energy exchange between the sites
\cite{franken2011,hannaford2002} is needed for delayed operation.
\item Passivity of the controller port together with a passive operator and environment yields stability of the interconnection by the passivity theorem; neither passivity is verified here.
\end{itemize}

%% =====================================================================
\section{Metrics}\label{sec:metrics}
%% =====================================================================
For errors $e[k,j]$ over $K$ samples and $J$ joints, $\mathrm{RMSE}_j=\sqrt{\tfrac1K\sum_ke[k,j]^2}$ and the \emph{pooled} RMSE is $\mathrm{RMSE}_{\rm pool}=\sqrt{\tfrac1{KJ}\sum_{k,j}e[k,j]^2}=\sqrt{\tfrac1J\sum_j\mathrm{RMSE}_j^2}$.

\begin{proposition}\label{prop:rmse}
$\tfrac1J\sum_j\mathrm{RMSE}_j\le\mathrm{RMSE}_{\rm pool}$, with equality iff all $\mathrm{RMSE}_j$ are equal.
\end{proposition}
\begin{proof}
By Jensen's inequality (convexity of $x\mapsto x^2$), $\big(\tfrac1J\sum_ja_j\big)^2\le\tfrac1J\sum_ja_j^2$ for $a_j=\mathrm{RMSE}_j$; equality iff all $a_j$ coincide.
\end{proof}
For the draft's Table~I the per-joint values $0.2250,0.2715,0.3609,0.3152,0.5366$ give a mean of $0.3418$ and a pooled value of $0.3583$ rad, so the table's ``Overall'' row is the pooled one and the $0.3418$ quoted in the conclusion was the mean of the joint RMSEs.
\test{test\_mean\_of\_joint\_rmse\_never\_exceeds\_pooled\_rmse, test\_pooled\_rmse\_is\_not\_mean\_of\_rmse}

The \emph{lag} of joint $j$ is $s^\star\Delta t$ with $s^\star=\arg\max_{0\le s\le s_{\max}}\rho(q_m[0{:}K{-}s],\,q_s[s{:}K])$ ($\rho$ the Pearson correlation), at a resolution of one sample (20\,ms); the \emph{lag-compensated RMSE} is the RMSE of $q_s[s^\star{:}]-q_m[{:}K{-}s^\star]$.
Definitions of the force metrics are in \texttt{docs/EVALUATION\_METHODOLOGY.md}.

%% =====================================================================
\section{From the draft's equations to the revised ones}\label{sec:map-eq}
%% =====================================================================
\begin{center}\small
\begin{tabular}{@{}p{1.6cm}p{5.2cm}p{7.2cm}@{}}
\toprule
Draft & Issue & Revised\\
\midrule
Eq.\,(1) & $\tau_m=\mathcal F\{k_t\odot I\}$ with $I$ in mA treated as A ($1000\times$ too large); two-motor joints counted once; nominal $\Delta t$ in the finite difference & $\tau_m=n_\mu\odot k_t\odot10^{-3}I$ (Assumption~\ref{as:act}); message time stamps in the velocity\\
Eq.\,(2) & $\tau_m-(C-G)$: sign of $G$, $C$ used as a vector, $M\ddot q$ ignored (Lemma~\ref{lem:static}) & slave: momentum observer, Theorem~\ref{thm:obs}; master: no estimate (it only echoes its own command)\\
Eq.\,(3),(4) & $\tau_{m,cmd}$ contains $\hat\tau_{ext,m}$: integrates the command, Lemma~\ref{lem:pole} & Eq.~\eqref{eq:master}\\
Eq.\,(5) & Euler step with the physical inertia as $M_{sd}$ and $\varepsilon=10^{-4}$ & design $M_d,D_d,K_d$ with the exact discretisation (Theorem~\ref{thm:zoh}); the operator-force term is removed with the master estimate\\
Eq.\,(6) & joint-by-joint reflection; called four-channel & Jacobian map \eqref{eq:force} and energy tank; with the master estimate removed the scheme has \emph{three} channels\\
Eq.\,(7) & torque MPC applied as $\tau_0^\star$; the code optimised a dead variable and ignored $q$ limits & position-command QP \eqref{eq:mpc}, offset-free, applied as $u_0^\star$\\
\bottomrule
\end{tabular}
\end{center}

%% =====================================================================
\section{What is not proven, and what hardware must confirm}\label{sec:limits}
%% =====================================================================
\begin{enumerate}[leftmargin=*,itemsep=2pt]
\item \textbf{Closed-loop stability} of the MPC (no terminal ingredients) and of the full teleoperation loop, in particular with delay.
\item \textbf{Model validity on the real arm}: decoupled second-order position loop with a constant disturbance (Assumption~\ref{as:loop}); $\omega_i,\zeta_i$ in the repository were identified on the \emph{simulated} slave.
Backlash, the Dynamixel velocity/acceleration profile and current limiting are not modelled.
\item \textbf{Torque constants and inertial parameters} (Remark~\ref{rem:bias}): the observer's accuracy is limited by $k_t$, link masses and friction, which must be calibrated.
\item \textbf{Energy accounting} is sampled-data (Section~\ref{sec:tank}); the damping-dissipation term $P_{in}$ assumes the master's and slave's damping terms are implemented as written.
\item The \textbf{dead zones, gains and $N$} were chosen in simulation.
\end{enumerate}

\begin{thebibliography}{99}\small
\bibitem{vanloan1978} C.~F. Van Loan, ``Computing integrals involving the matrix exponential,'' \emph{IEEE Trans. Autom. Control}, vol.~23, no.~3, 1978, doi:10.1109/TAC.1978.1101743.
\bibitem{delucamattone2003} A.~De~Luca and R.~Mattone, ``Actuator failure detection and isolation using generalized momenta,'' in \emph{Proc. IEEE ICRA}, 2003, vol.~1, pp.~634--639.
\bibitem{haddadin2017} S.~Haddadin, A.~De~Luca, and A.~Albu-Sch\"affer, ``Robot collisions: A survey on detection, isolation, and identification,'' \emph{IEEE Trans. Robot.}, vol.~33, no.~6, pp.~1292--1312, 2017.
\bibitem{stellato2020osqp} B.~Stellato, G.~Banjac, P.~Goulart, A.~Bemporad, and S.~Boyd, ``OSQP: an operator splitting solver for quadratic programs,'' \emph{Math. Program. Comput.}, vol.~12, no.~4, pp.~637--672, 2020.
\bibitem{boyd2011admm} S.~Boyd, N.~Parikh, E.~Chu, B.~Peleato, and J.~Eckstein, ``Distributed optimization and statistical learning via the alternating direction method of multipliers,'' \emph{Found. Trends Mach. Learn.}, vol.~3, no.~1, pp.~1--122, 2011.
\bibitem{muske2002} K.~R. Muske and T.~A. Badgwell, ``Disturbance modeling for offset-free linear model predictive control,'' \emph{J. Process Control}, vol.~12, no.~5, pp.~617--632, 2002.
\bibitem{pannocchia2003} G.~Pannocchia and J.~B. Rawlings, ``Disturbance models for offset-free model-predictive control,'' \emph{AIChE J.}, vol.~49, no.~2, pp.~426--437, 2003.
\bibitem{franken2011} M.~Franken, S.~Stramigioli, S.~Misra, C.~Secchi, and A.~Macchelli, ``Bilateral telemanipulation with time delays: A two-layer approach combining passivity and transparency,'' \emph{IEEE Trans. Robot.}, vol.~27, no.~4, pp.~741--756, 2011.
\bibitem{hannaford2002} B.~Hannaford and J.-H. Ryu, ``Time-domain passivity control of haptic interfaces,'' \emph{IEEE Trans. Robot. Autom.}, vol.~18, no.~1, pp.~1--10, 2002.
\bibitem{anderson1989} R.~J. Anderson and M.~W. Spong, ``Bilateral control of teleoperators with time delay,'' \emph{IEEE Trans. Autom. Control}, vol.~34, no.~5, pp.~494--501, 1989.
\bibitem{lawrence1993} D.~A. Lawrence, ``Stability and transparency in bilateral teleoperation,'' \emph{IEEE Trans. Robot. Autom.}, vol.~9, no.~5, pp.~624--637, 1993.
\end{thebibliography}
\end{document}
